\begin{table}[!t]
\centering
\footnotesize
\setlength{\tabcolsep}{2.5pt}
\caption{Primary results under controlled reference sampling. EPA and VitalDB use their localized complete-data empirical benchmarks; AI uses its population benchmark.}
\label{tab:application-primary}
\begin{tabular}{llrrrrrl}
\toprule
Study & Target & Ref. frac. & Anchor & Full & Adaptive & Mean wt. & Adaptive/Anchor [95\% MC] \\
\midrule
EPA & 30\% RH & 10.0\% & 2.060 & 0.738 & 0.818 & 0.846 & 0.397 [0.337, 0.457] \\
 & 50\% RH & 10.0\% & 0.864 & 0.375 & 0.386 & 0.888 & 0.447 [0.385, 0.510] \\
 & 70\% RH & 10.0\% & 1.601 & 0.722 & 0.765 & 0.862 & 0.478 [0.423, 0.534] \\
VitalDB & 30 min & 10.0\% & 0.205 & 0.177 & 0.181 & 0.803 & 0.882 [0.851, 0.914] \\
 & 75 min & 10.0\% & 0.393 & 0.347 & 0.349 & 0.867 & 0.888 [0.857, 0.919] \\
 & 180 min & 10.0\% & 0.241 & 0.232 & 0.234 & 0.672 & 0.971 [0.949, 0.992] \\
AI & Random review & 2.5\% & 9.091 & 3.380 & 4.195 & 0.775 & 0.461 [0.427, 0.496] \\
 & Random review & 10.0\% & 2.147 & 0.814 & 0.861 & 0.891 & 0.401 [0.363, 0.439] \\
 & Random review & 40.0\% & 0.434 & 0.197 & 0.199 & 0.955 & 0.457 [0.421, 0.494] \\
\bottomrule
\end{tabular}
\vspace{2pt}
\parbox{0.98\textwidth}{\footnotesize ISE values are multiplied by $10^3$. Risk-ratio intervals use common reference samples and quantify Monte Carlo variation conditional on the stated study-specific benchmark; they are not population-sampling confidence intervals.}
\end{table}
